\documentclass[11pt]{article}

\usepackage[utf8]{inputenc}
\usepackage[T1]{fontenc}
\usepackage[spanish,es-nodecimaldot]{babel}
\usepackage{lmodern}
\usepackage[a4paper,margin=2.2cm]{geometry}
\usepackage{microtype}
\usepackage{amsmath,amssymb,mathtools}
\usepackage{booktabs,array,multirow}
\usepackage{enumitem}
\usepackage{xcolor}
\usepackage{hyperref}
\usepackage{tcolorbox}
\usepackage{tikz}
\usetikzlibrary{arrows.meta,positioning,shapes.geometric,calc,fit,matrix}
\usepackage{pgfplots}
\pgfplotsset{compat=1.18}
\usepackage{listings}
\usepackage{fancyhdr}
\usepackage{lastpage}

\definecolor{azul}{HTML}{1F4E79}
\definecolor{azulclaro}{HTML}{EAF3F8}
\definecolor{verde}{HTML}{2D6A4F}
\definecolor{verdeclaro}{HTML}{E9F5EE}
\definecolor{gris}{HTML}{555555}
\definecolor{rojoclaro}{HTML}{FBEAEA}

\hypersetup{colorlinks=true,linkcolor=azul,urlcolor=azul,citecolor=azul}
\setlist[itemize]{leftmargin=1.6em,itemsep=2pt,topsep=4pt}
\setlist[enumerate]{leftmargin=2em,itemsep=4pt,topsep=4pt}
\setlength{\parindent}{0pt}
\setlength{\parskip}{5pt}
\setlength{\headheight}{14pt}
\pagestyle{fancy}
\fancyhf{}
\lhead{Aprendizaje de Máquina -- Parcial I}
\rhead{Guía resuelta}
\cfoot{\thepage\ de \pageref{LastPage}}

\newtcolorbox{respuesta}{
  colback=azulclaro,
  colframe=azul,
  boxrule=0.7pt,
  arc=2mm,
  left=2mm,right=2mm,top=1mm,bottom=1mm
}
\newtcolorbox{clave}{
  colback=verdeclaro,
  colframe=verde,
  boxrule=0.7pt,
  arc=2mm,
  title=Idea clave,
  fonttitle=\bfseries
}
\newtcolorbox{advertencia}{
  colback=rojoclaro,
  colframe=red!70!black,
  boxrule=0.7pt,
  arc=2mm,
  title=Atención,
  fonttitle=\bfseries
}

\newcommand{\pregunta}[1]{\vspace{2mm}\textbf{#1}}
\newcommand{\R}{\mathbb{R}}
\newcommand{\yhat}{\hat y}

\lstset{
  language=Python,
  basicstyle=\ttfamily\small,
  backgroundcolor=\color{black!3},
  frame=single,
  rulecolor=\color{black!20},
  keywordstyle=\color{azul}\bfseries,
  commentstyle=\color{verde},
  showstringspaces=false,
  columns=fullflexible,
  breaklines=true
}

\title{\textbf{Guía resuelta -- Examen Parcial I}\\[2mm]
\large Aprendizaje de Máquina -- LDS1071}
\author{Material de estudio}
\date{}

\begin{document}
\maketitle
\tableofcontents
\newpage

\section*{Notación y vocabulario mínimo}
\addcontentsline{toc}{section}{Notación y vocabulario mínimo}

Para un conjunto de $n$ observaciones y $p$ características, usaremos
\[
X\in\R^{n\times p},\qquad y\in\R^n\ \text{o}\ y\in\{1,\dots,K\}^n.
\]
La fila $x_i^\top$ representa la observación $i$; la columna $j$ representa la característica $j$; $y_i$ es la variable objetivo de la observación $i$; y $\hat y_i$ es la predicción del modelo.

\begin{clave}
El hilo conductor de casi todo el parcial es: \emph{aprender patrones con datos disponibles y medir si esos patrones funcionan en datos no vistos}. Entrenar bien no basta; hay que generalizar.
\end{clave}

% ============================================================
\section{Fundamentos del aprendizaje de máquina}

\subsection*{Mapa conceptual}
\begin{center}
\begin{tikzpicture}[
  node distance=10mm and 11mm,
  box/.style={draw,rounded corners,align=center,minimum width=3.0cm,minimum height=9mm,fill=blue!4},
  arr/.style={-Latex,thick}
]
\node[box] (ml) {Aprendizaje\\de máquina};
\node[box,below left=of ml] (sup) {Supervisado\\hay objetivo $y$};
\node[box,below right=of ml] (unsup) {No supervisado\\no hay objetivo $y$};
\node[box,below left=of sup] (class) {Clasificación\\$y$ categórica};
\node[box,below right=of sup] (reg) {Regresión\\$y$ numérica};
\node[box,below=of unsup] (clus) {Clustering\\descubrir grupos};
\draw[arr] (ml)--(sup); \draw[arr] (ml)--(unsup);
\draw[arr] (sup)--(class); \draw[arr] (sup)--(reg); \draw[arr] (unsup)--(clus);
\end{tikzpicture}
\end{center}

\pregunta{1. ¿Qué diferencia fundamental existe entre aprendizaje supervisado y no supervisado?}
\begin{respuesta}
En el aprendizaje \textbf{supervisado} cada observación de entrenamiento incluye una salida conocida $y_i$ y se busca aprender una función $f(X)\approx y$. En el aprendizaje \textbf{no supervisado} no existe una etiqueta objetivo conocida; el algoritmo intenta descubrir estructura en $X$, por ejemplo grupos, direcciones principales o representaciones latentes.
\end{respuesta}

\pregunta{2. ¿Qué información adicional tienen los datos utilizados en aprendizaje supervisado?}
\begin{respuesta}
Tienen una \textbf{etiqueta o variable objetivo} asociada a cada observación. Matemáticamente se observan pares $(x_i,y_i)$, no sólo $x_i$.
\end{respuesta}

\pregunta{3. ¿Cuál es la diferencia entre clasificación y regresión?}
\begin{respuesta}
Ambos son aprendizaje supervisado. En \textbf{clasificación}, $y$ toma categorías discretas, como «fraude/no fraude». En \textbf{regresión}, $y$ es una cantidad numérica, como precio, temperatura o gasto mensual.
\end{respuesta}

\pregunta{4. ¿Puede un problema de clasificación tener más de dos clases?}
\begin{respuesta}
Sí. Si $K>2$ se llama \textbf{clasificación multiclase}. Por ejemplo, clasificar una imagen como gato, perro, caballo o ave tiene cuatro clases.
\end{respuesta}

\pregunta{5. ¿Qué diferencia existe entre clasificación binaria y multiclase?}
\begin{respuesta}
La clasificación binaria tiene exactamente dos clases, típicamente codificadas como $0$ y $1$. La multiclase tiene tres o más categorías posibles. El objetivo sigue siendo asignar una categoría a cada observación.
\end{respuesta}

\pregunta{6. ¿Qué significa que un conjunto de datos esté etiquetado?}
\begin{respuesta}
Que para cada observación de entrada se conoce la salida correcta o variable objetivo. Por ejemplo, si cada correo viene acompañado de la etiqueta «spam» o «no spam», los datos están etiquetados.
\end{respuesta}

\pregunta{7. ¿Qué tipo de aprendizaje utilizarías para descubrir grupos naturales dentro de un conjunto de datos?}
\begin{respuesta}
Aprendizaje \textbf{no supervisado}, específicamente un método de \textbf{clustering} como $k$-means, porque no conocemos de antemano las etiquetas de los grupos.
\end{respuesta}

\pregunta{8. ¿Qué tipo de aprendizaje utilizarías para predecir una cantidad continua?}
\begin{respuesta}
Aprendizaje supervisado de \textbf{regresión}. Ejemplo: predecir el precio de una vivienda en pesos.
\end{respuesta}

\pregunta{9. Da dos ejemplos propios de problemas de clasificación.}
\begin{respuesta}
(1) Determinar si una transacción bancaria es fraudulenta o legítima. (2) Clasificar una fotografía de una planta según su especie.
\end{respuesta}

\pregunta{10. Da dos ejemplos propios de problemas de regresión.}
\begin{respuesta}
(1) Predecir el consumo eléctrico de una casa durante el siguiente día. (2) Predecir el rendimiento de un cultivo en toneladas por hectárea.
\end{respuesta}

\pregunta{11. Da dos ejemplos propios de problemas de aprendizaje no supervisado.}
\begin{respuesta}
(1) Segmentar clientes en grupos con comportamientos de compra semejantes sin categorías previas. (2) Detectar estructuras o agrupaciones en documentos usando sólo sus características textuales.
\end{respuesta}

% ============================================================
\section{Datos: características y variable objetivo}

\subsection*{Representación matricial}
\begin{center}
\begin{tikzpicture}[scale=0.92]
\matrix[matrix of math nodes,left delimiter={[},right delimiter={]},nodes={minimum width=9mm,minimum height=6mm}] (X) {
 x_{11} & x_{12} & \cdots & x_{1p}\\
 x_{21} & x_{22} & \cdots & x_{2p}\\
 \vdots & \vdots & \ddots & \vdots\\
 x_{n1} & x_{n2} & \cdots & x_{np}\\
};
\node[above=2mm of X] {$X$: entradas / características};
\draw[decorate,decoration={brace,amplitude=4pt},thick] ($(X.north west)+(-.4,.1)$) -- ($(X.south west)+(-.4,-.1)$) node[midway,left=7mm] {$n$ observaciones};
\draw[decorate,decoration={brace,amplitude=4pt,mirror},thick] ($(X.south west)+(0,-.35)$) -- ($(X.south east)+(0,-.35)$) node[midway,below=6mm] {$p$ características};
\node[right=18mm of X] (arrow) {$\longrightarrow$};
\matrix[matrix of math nodes,left delimiter={[},right delimiter={]},right=10mm of arrow] (y) {y_1\\y_2\\\vdots\\y_n\\};
\node[above=2mm of y] {$y$: objetivo};
\end{tikzpicture}
\end{center}

\pregunta{1. ¿Qué representa normalmente una fila de la matriz $X$?}
\begin{respuesta}
Una \textbf{observación, instancia o ejemplo}. Por ejemplo, una fila puede corresponder a un cliente, paciente, casa o imagen.
\end{respuesta}

\pregunta{2. ¿Qué representa normalmente una columna de $X$?}
\begin{respuesta}
Una \textbf{característica} o variable de entrada medida para todas las observaciones: edad, ingreso, área, temperatura, etc.
\end{respuesta}

\pregunta{3. ¿Qué representa $y$ en aprendizaje supervisado?}
\begin{respuesta}
La \textbf{variable objetivo} que el modelo intenta predecir. En clasificación contiene clases; en regresión contiene valores numéricos.
\end{respuesta}

\pregunta{4. Si tenemos 1000 observaciones y 15 características, ¿qué dimensiones tendría normalmente $X$?}
\begin{respuesta}
$X$ tendría dimensión $1000\times 15$: 1000 filas y 15 columnas.
\end{respuesta}

\pregunta{5. ¿Cuántos valores tendría $y$ en el ejemplo anterior?}
\begin{respuesta}
1000 valores, uno por observación. Puede escribirse $y\in\R^{1000}$ en regresión.
\end{respuesta}

\pregunta{6. ¿Puede una característica ser categórica?}
\begin{respuesta}
Sí. Por ejemplo: ciudad, tipo de cliente o color. Algunos modelos admiten categorías directamente; otros requieren una codificación numérica como \emph{one-hot encoding}.
\end{respuesta}

\pregunta{7. ¿Puede una característica ser numérica?}
\begin{respuesta}
Sí. Ejemplos: edad, ingreso, estatura, número de compras o temperatura.
\end{respuesta}

\pregunta{8. ¿Qué diferencia existe entre una característica y la variable objetivo?}
\begin{respuesta}
Una característica forma parte de la información de entrada que el modelo puede utilizar para producir una predicción. La variable objetivo es aquello que queremos predecir y proporciona la señal de aprendizaje durante el entrenamiento supervisado.
\end{respuesta}

\pregunta{9. Dado un problema real, identifica posibles características y la variable objetivo.}
\begin{respuesta}
Ejemplo: predecir el precio de una vivienda. Características: superficie, número de habitaciones, antigüedad, ubicación y distancia al centro. Variable objetivo: precio de venta.
\end{respuesta}

\pregunta{10. ¿Por qué la variable objetivo no se considera normalmente una característica de entrada?}
\begin{respuesta}
Porque en el momento real de predicción precisamente \textbf{no conocemos} esa variable. Incluirla como entrada produciría \textbf{fuga de información} (\emph{data leakage}): el modelo recibiría parte o toda la respuesta que debería inferir.
\end{respuesta}

% ============================================================
\section{Entrenamiento, validación y prueba}

\begin{center}
\begin{tikzpicture}[x=1cm,y=1cm]
\draw[rounded corners,fill=blue!8] (0,0) rectangle (12,1);
\fill[blue!35] (0,0) rectangle (8.4,1);
\fill[green!35] (8.4,0) rectangle (10.2,1);
\fill[orange!45] (10.2,0) rectangle (12,1);
\node at (4.2,.5) {Entrenamiento $\sim70\%$};
\node at (9.3,.5) {Validación};
\node at (11.1,.5) {Prueba};
\node[below] at (4.2,-.15) {ajustar parámetros};
\node[below] at (9.3,-.15) {elegir modelo};
\node[below] at (11.1,-.15) {evaluación final};
\end{tikzpicture}
\end{center}

\pregunta{1. ¿Sobre qué conjunto se ajustan normalmente los parámetros del modelo?}
\begin{respuesta}
Sobre el \textbf{conjunto de entrenamiento}. Los parámetros son los valores internos aprendidos por el algoritmo a partir de $(X_{train},y_{train})$.
\end{respuesta}

\pregunta{2. ¿Para qué se utiliza el conjunto de validación?}
\begin{respuesta}
Para comparar modelos, seleccionar hiperparámetros, decidir complejidad y tomar decisiones de desarrollo sin usar el conjunto de prueba final.
\end{respuesta}

\pregunta{3. ¿Para qué se utiliza el conjunto de prueba?}
\begin{respuesta}
Para estimar, al final del desarrollo, el desempeño del modelo sobre datos no vistos. Debe funcionar como una evaluación final lo más independiente posible.
\end{respuesta}

\pregunta{4. ¿Por qué no deberíamos entrenar y evaluar un modelo utilizando exactamente los mismos datos?}
\begin{respuesta}
Porque mediríamos qué tan bien el modelo reproduce datos que ya vio, no qué tan bien generaliza. Un modelo puede memorizar peculiaridades del entrenamiento y aun así fallar en datos nuevos.
\end{respuesta}

\pregunta{5. ¿Qué significa que un modelo generalice?}
\begin{respuesta}
Que mantenga un buen desempeño en observaciones nuevas provenientes del mismo fenómeno o distribución de interés, no sólo en los ejemplos usados para ajustarlo.
\end{respuesta}

\pregunta{6. ¿Qué problema puede surgir si seleccionamos nuestro modelo observando continuamente los resultados del conjunto de prueba?}
\begin{respuesta}
El desarrollo termina adaptándose indirectamente al conjunto de prueba. Aunque no entrenemos formalmente con él, nuestras decisiones usan su información; por ello la estimación final puede quedar \textbf{optimistamente sesgada}.
\end{respuesta}

\pregunta{7. ¿Qué diferencia existe entre seleccionar un modelo y evaluar el modelo final?}
\begin{respuesta}
Seleccionar un modelo es comparar alternativas y tomar decisiones, normalmente con validación. Evaluar el modelo final consiste en estimar su capacidad de generalización una vez concluidas esas decisiones, usando datos de prueba que no guiaron el desarrollo.
\end{respuesta}

\pregunta{8. Si tienes 10,000 observaciones, propón una división razonable entre entrenamiento, validación y prueba.}
\begin{respuesta}
Una división razonable es 70\% / 15\% / 15\%: 7000 para entrenamiento, 1500 para validación y 1500 para prueba. También pueden ser razonables otras proporciones según tamaño, ruido y costo de obtener datos.
\end{respuesta}

\pregunta{9. ¿Existe una única proporción correcta para dividir los datos?}
\begin{respuesta}
No. La división depende de cuántos datos haya, cuánta variabilidad exista, el costo de entrenar, la necesidad de precisión en la evaluación y si se utiliza validación cruzada.
\end{respuesta}

\pregunta{10. ¿Por qué normalmente queremos utilizar una mayor cantidad de datos para entrenamiento que para prueba?}
\begin{respuesta}
Porque los datos de entrenamiento contienen la información con la que se aprenden los patrones del modelo. El conjunto de prueba sólo necesita ser suficientemente grande y representativo para estimar el error con precisión razonable.
\end{respuesta}

% ============================================================
\section{Generalización}

\pregunta{1. ¿Qué significa que un modelo tenga buena capacidad de generalización?}
\begin{respuesta}
Que su error en datos no vistos sea bajo y semejante al esperado para observaciones futuras del problema real.
\end{respuesta}

\pregunta{2. ¿Un error de entrenamiento muy pequeño garantiza un buen modelo? Explica.}
\begin{respuesta}
No. Puede indicar un ajuste correcto, pero también memorización. El criterio relevante es el desempeño fuera de muestra: un error de entrenamiento pequeño con error de prueba grande es una señal típica de sobreajuste.
\end{respuesta}

\pregunta{3. ¿Qué información nos proporciona comparar el desempeño en entrenamiento y en datos no observados?}
\begin{respuesta}
Permite separar dos preguntas: (i) si el modelo logra ajustarse a los patrones disponibles y (ii) si esos patrones se transfieren a datos nuevos. La brecha entre errores ayuda a diagnosticar sobreajuste; errores altos en ambos conjuntos sugieren subajuste.
\end{respuesta}

\pregunta{4. Si dos modelos tienen el mismo error de entrenamiento pero diferente error de prueba, ¿qué información puedes obtener?}
\begin{respuesta}
El modelo con menor error de prueba muestra mejor generalización en esa evaluación. Como ambos ajustan igual el entrenamiento, la diferencia aparece al enfrentarse a datos no vistos.
\end{respuesta}

\pregunta{5. ¿Por qué el desempeño sobre datos no observados es importante en aplicaciones reales?}
\begin{respuesta}
Porque el propósito de entrenar un modelo es usarlo después sobre casos que todavía no conocemos. Un sistema que sólo funciona en los datos históricos no cumple el objetivo predictivo.
\end{respuesta}

% ============================================================
\section{Underfitting y overfitting}

\subsection*{Curva conceptual error--complejidad}
\begin{center}
\begin{tikzpicture}
\begin{axis}[
  width=11cm,height=6.2cm,
  xlabel={Complejidad del modelo}, ylabel={Error},
  xmin=0,xmax=10,ymin=0,ymax=8,
  xtick=\empty,ytick=\empty,
  axis lines=left,
  legend style={at={(0.5,-0.18)},anchor=north,legend columns=2,draw=none}
]
\addplot[thick,blue,domain=0.3:9.7,samples=100]{6.2*exp(-0.35*x)+0.45};
\addlegendentry{Entrenamiento}
\addplot[thick,red,domain=0.3:9.7,samples=100]{0.16*(x-5.2)^2+1.6};
\addlegendentry{Validación/prueba}
\draw[dashed] (axis cs:5.2,0) -- (axis cs:5.2,2.0);
\node at (axis cs:1.5,7.2) {Subajuste};
\node at (axis cs:8.5,7.2) {Sobreajuste};
\node at (axis cs:5.2,2.45) {zona útil};
\end{axis}
\end{tikzpicture}
\end{center}

\pregunta{1. ¿Cómo esperarías que fueran los errores de entrenamiento y prueba cuando existe subajuste?}
\begin{respuesta}
Ambos suelen ser relativamente altos y cercanos. El modelo es demasiado simple o insuficiente incluso para capturar los patrones del entrenamiento.
\end{respuesta}

\pregunta{2. ¿Cómo esperarías que fueran cuando existe sobreajuste?}
\begin{respuesta}
El error de entrenamiento suele ser muy bajo, mientras que el error de validación o prueba es notablemente mayor. La \textbf{brecha de generalización} es grande.
\end{respuesta}

\pregunta{3. ¿Por qué aumentar excesivamente la complejidad de un modelo puede favorecer el sobreajuste?}
\begin{respuesta}
Porque aumenta la capacidad de representar no sólo patrones estables sino también ruido, anomalías y coincidencias particulares del conjunto de entrenamiento.
\end{respuesta}

\pregunta{4. ¿Por qué un modelo demasiado simple puede presentar subajuste?}
\begin{respuesta}
Porque su familia de funciones no tiene suficiente flexibilidad para aproximar la relación real entre entradas y objetivo. El error permanece alto aun en entrenamiento.
\end{respuesta}

\pregunta{5. Menciona dos posibles estrategias para combatir el sobreajuste.}
\begin{respuesta}
(1) Reducir la complejidad o aplicar regularización. (2) Obtener más datos de entrenamiento representativos. También pueden ayudar selección de características, \emph{early stopping} y validación adecuada.
\end{respuesta}

\pregunta{6. Menciona dos posibles estrategias para combatir el subajuste.}
\begin{respuesta}
(1) Usar un modelo más flexible o aumentar su complejidad. (2) Construir características más informativas. También puede ser necesario entrenar mejor o reducir una regularización excesiva.
\end{respuesta}

\pregunta{7. ¿Obtener más datos de entrenamiento puede ayudar contra el sobreajuste?}
\begin{respuesta}
Sí, frecuentemente. Más datos dificultan memorizar peculiaridades individuales y permiten estimar patrones más estables, aunque no corrigen por sí solos todos los problemas de diseño o sesgo.
\end{respuesta}

\pregunta{8. ¿Reducir la complejidad del modelo puede ayudar contra el sobreajuste?}
\begin{respuesta}
Sí. Limitar la flexibilidad reduce la capacidad de ajustar ruido. Debe hacerse sin simplificar tanto que aparezca subajuste.
\end{respuesta}

\pregunta{9. ¿Aumentar la complejidad siempre mejora el desempeño?}
\begin{respuesta}
No. Normalmente reduce el error de entrenamiento, pero el error fuera de muestra puede disminuir sólo hasta cierto punto y luego aumentar debido al sobreajuste.
\end{respuesta}

\pregunta{10. ¿Qué buscarías idealmente al comparar los errores de entrenamiento y prueba?}
\begin{respuesta}
Errores bajos en ambos conjuntos y una brecha razonablemente pequeña. No buscamos igualdad exacta, sino buen desempeño fuera de muestra sin señales fuertes de incapacidad o memorización.
\end{respuesta}

% ============================================================
\section{Evaluación de modelos de regresión}

Para errores $e_i=y_i-\hat y_i$:
\[
\mathrm{MAE}=\frac1n\sum_{i=1}^n |e_i|,\qquad
\mathrm{MSE}=\frac1n\sum_{i=1}^n e_i^2,\qquad
\mathrm{RMSE}=\sqrt{\mathrm{MSE}}.
\]

\pregunta{1. ¿Qué mide el MAE?}
\begin{respuesta}
El tamaño promedio del error en valor absoluto. Puede interpretarse como cuánto se equivoca el modelo, en promedio, ignorando el signo del error.
\end{respuesta}

\pregunta{2. ¿Qué mide el MSE?}
\begin{respuesta}
El promedio de los errores al cuadrado. Penaliza de forma creciente los errores grandes y es siempre no negativo.
\end{respuesta}

\pregunta{3. ¿Por qué los errores grandes tienen mayor influencia en el MSE?}
\begin{respuesta}
Porque el error se eleva al cuadrado: duplicar un error multiplica su contribución por cuatro. Por ejemplo, un error de $10$ aporta $100$, mientras uno de $2$ aporta sólo $4$.
\end{respuesta}

\pregunta{4. ¿Qué ventaja de interpretación tiene RMSE respecto a MSE?}
\begin{respuesta}
El RMSE vuelve a las mismas unidades de la variable objetivo. Si $y$ está en pesos, RMSE está en pesos; MSE está en pesos cuadrados.
\end{respuesta}

\pregunta{5. Si la variable objetivo está medida en pesos, ¿en qué unidades se expresa MAE?}
\begin{respuesta}
En pesos.
\end{respuesta}

\pregunta{6. ¿En qué unidades se expresa RMSE?}
\begin{respuesta}
En las mismas unidades de $y$; por ejemplo, pesos.
\end{respuesta}

\pregunta{7. ¿Un valor más grande o más pequeño de MAE representa mejores predicciones?}
\begin{respuesta}
Más pequeño. El valor ideal es $0$, que significaría predicciones exactas en todas las observaciones evaluadas.
\end{respuesta}

\pregunta{8. Calcula MAE para $y=(10,20,30)$, $\hat y=(12,18,27)$.}
\begin{respuesta}
Los errores son $(-2,2,3)$ y los errores absolutos $(2,2,3)$. Por tanto,
\[
\mathrm{MAE}=\frac{2+2+3}{3}=\frac73\approx 2.333.
\]
\end{respuesta}

\pregunta{9. Calcula MSE para los mismos datos.}
\begin{respuesta}
\[
\mathrm{MSE}=\frac{(-2)^2+2^2+3^2}{3}=\frac{4+4+9}{3}=\frac{17}{3}\approx5.667.
\]
\end{respuesta}

\pregunta{10. Calcula RMSE para los mismos datos.}
\begin{respuesta}
\[
\mathrm{RMSE}=\sqrt{\frac{17}{3}}\approx 2.380.
\]
\end{respuesta}

% ============================================================
\section{Evaluación de modelos de clasificación}

\subsection*{Matriz de confusión binaria}
\begin{center}
\begin{tikzpicture}
\matrix[matrix of nodes,
  nodes={draw,minimum width=3cm,minimum height=1.2cm,align=center},
  row 1/.style={nodes={fill=black!6,font=\bfseries}},
  column 1/.style={nodes={fill=black!6,font=\bfseries}}
] (m) {
 & Predicho $+$ & Predicho $-$ \\
Real $+$ & TP & FN \\
Real $-$ & FP & TN \\
};
\node[above=4mm of m] {Las filas representan la realidad; las columnas, la predicción};
\end{tikzpicture}
\end{center}

\[
\mathrm{Accuracy}=\frac{TP+TN}{TP+TN+FP+FN},\quad
\mathrm{Precision}=\frac{TP}{TP+FP},\quad
\mathrm{Recall}=\frac{TP}{TP+FN}.
\]
\[
F_1=2\frac{\mathrm{Precision}\cdot\mathrm{Recall}}{\mathrm{Precision}+\mathrm{Recall}}.
\]

\pregunta{1. ¿Qué significa un falso positivo?}
\begin{respuesta}
El modelo predice la clase positiva, pero la observación en realidad pertenece a la clase negativa. Es una «alarma» incorrecta.
\end{respuesta}

\pregunta{2. ¿Qué significa un falso negativo?}
\begin{respuesta}
El modelo predice la clase negativa, pero la observación era realmente positiva. Es un caso positivo que el sistema no detectó.
\end{respuesta}

\pregunta{3. ¿Qué mide accuracy?}
\begin{respuesta}
La fracción total de predicciones correctas: positivos y negativos correctamente clasificados sobre el total de observaciones.
\end{respuesta}

\pregunta{4. ¿Qué mide precision?}
\begin{respuesta}
Entre todo lo que el modelo \textbf{predijo como positivo}, qué proporción realmente era positiva. Responde: «cuando el modelo dice positivo, ¿qué tan frecuentemente tiene razón?»
\end{respuesta}

\pregunta{5. ¿Qué mide recall?}
\begin{respuesta}
Entre todos los casos \textbf{realmente positivos}, qué proporción fue detectada. Responde: «de los positivos existentes, ¿cuántos encontró el modelo?»
\end{respuesta}

\pregunta{6. ¿Qué combina el F1-score?}
\begin{respuesta}
Combina precision y recall mediante su media armónica, por lo que un valor alto requiere que ambas sean razonablemente altas.
\end{respuesta}

\pregunta{7. Construye una matriz de confusión a partir de una lista pequeña de predicciones.}
\begin{respuesta}
Ejemplo: realidad $(1,0,1,1,0,0)$ y predicción $(1,1,1,0,0,0)$. Se obtiene $TP=2$, $FN=1$, $FP=1$, $TN=2$:
\[
\begin{array}{c|cc}
&\hat y=1&\hat y=0\\\hline
y=1&2&1\\
y=0&1&2
\end{array}
\]
\end{respuesta}

\pregunta{8. Calcula accuracy a partir de una matriz de confusión.}
\begin{respuesta}
Se suman los aciertos de la diagonal y se dividen entre el total: $\mathrm{Accuracy}=(TP+TN)/N$. En el ejemplo anterior: $(2+2)/6=2/3\approx66.7\%$.
\end{respuesta}

\pregunta{9. Calcula precision.}
\begin{respuesta}
$\mathrm{Precision}=TP/(TP+FP)$. En el ejemplo: $2/(2+1)=2/3\approx66.7\%$.
\end{respuesta}

\pregunta{10. Calcula recall.}
\begin{respuesta}
$\mathrm{Recall}=TP/(TP+FN)$. En el ejemplo: $2/(2+1)=2/3\approx66.7\%$.
\end{respuesta}

\pregunta{11. Calcula F1-score.}
\begin{respuesta}
Con precision $=2/3$ y recall $=2/3$,
\[
F_1=2\frac{(2/3)(2/3)}{2/3+2/3}=\frac23\approx66.7\%.
\]
\end{respuesta}

\pregunta{12. Da un ejemplo donde un falso positivo sea especialmente costoso.}
\begin{respuesta}
Un sistema de bloqueo automático que marque una transacción legítima como fraude y cancele una compra crítica. El costo proviene de actuar sobre una alarma incorrecta.
\end{respuesta}

\pregunta{13. Da un ejemplo donde un falso negativo sea especialmente costoso.}
\begin{respuesta}
Un detector de fraude que clasifique una transacción fraudulenta como legítima. El sistema deja pasar precisamente el caso que debía detectar.
\end{respuesta}

\pregunta{14. ¿Por qué la métrica apropiada depende del problema?}
\begin{respuesta}
Porque distintos errores tienen costos distintos. Si es muy grave perder positivos, interesa recall; si es muy grave generar falsas alarmas, interesa precision; si las clases están balanceadas y los costos son comparables, accuracy puede ser informativa.
\end{respuesta}

% ============================================================
\section{Clases desbalanceadas}

\pregunta{1. ¿Qué significa que las clases estén desbalanceadas?}
\begin{respuesta}
Que algunas clases contienen muchas más observaciones que otras. Por ejemplo, 98\% negativos y 2\% positivos.
\end{respuesta}

\pregunta{2. Si el 99\% de las observaciones pertenece a una misma clase, ¿qué accuracy obtiene un clasificador que siempre predice esa clase?}
\begin{respuesta}
99\% de accuracy, aunque no detecte ningún ejemplo de la clase minoritaria.
\end{respuesta}

\pregunta{3. ¿Por qué ese resultado puede ser engañoso?}
\begin{respuesta}
Porque el accuracy está dominado por la clase mayoritaria. El modelo puede parecer excelente y, al mismo tiempo, tener recall igual a cero para la clase minoritaria.
\end{respuesta}

\pregunta{4. ¿Qué otras métricas revisarías?}
\begin{respuesta}
Precision, recall, F1-score y la matriz de confusión. También puede ser útil medir métricas por clase, según el objetivo del problema.
\end{respuesta}

\pregunta{5. ¿Cuándo podría ser particularmente importante recall?}
\begin{respuesta}
Cuando perder un positivo real es muy costoso. Por ejemplo, en una primera etapa de detección de fraude interesa descubrir la mayor fracción posible de transacciones fraudulentas.
\end{respuesta}

\pregunta{6. ¿Cuándo podría ser particularmente importante precision?}
\begin{respuesta}
Cuando actuar sobre un falso positivo sea costoso. Por ejemplo, si cada alarma genera una revisión manual muy cara, queremos que una gran proporción de las alarmas sea correcta.
\end{respuesta}

\pregunta{7. ¿Por qué no existe una única métrica que sea siempre la mejor para cualquier problema?}
\begin{respuesta}
Porque las métricas enfatizan aspectos distintos del error y los costos del mundo real cambian entre aplicaciones. La elección de métrica forma parte de la definición del objetivo.
\end{respuesta}

% ============================================================
\section{Parámetros e hiperparámetros}

\begin{center}
\begin{tikzpicture}[node distance=12mm,box/.style={draw,rounded corners,align=center,minimum width=5cm,minimum height=1cm}]
\node[box,fill=blue!6] (p) {Parámetros\\aprendidos de los datos};
\node[box,fill=green!8,right=18mm of p] (h) {Hiperparámetros\\configuran el aprendizaje};
\node[below=7mm of p,align=center] {Ej.: coeficientes de regresión,\\pesos de una red};
\node[below=7mm of h,align=center] {Ej.: profundidad de árbol,\\fuerza de regularización};
\end{tikzpicture}
\end{center}

\pregunta{1. ¿Cuál es la diferencia entre parámetro e hiperparámetro?}
\begin{respuesta}
Un \textbf{parámetro} se estima a partir de los datos durante el entrenamiento. Un \textbf{hiperparámetro} controla el algoritmo, la arquitectura o la complejidad y se elige fuera del ajuste ordinario de parámetros.
\end{respuesta}

\pregunta{2. ¿Quién determina los parámetros del modelo?}
\begin{respuesta}
El algoritmo de entrenamiento, optimizando algún criterio con los datos de entrenamiento.
\end{respuesta}

\pregunta{3. ¿Cómo se pueden seleccionar los hiperparámetros?}
\begin{respuesta}
Comparando configuraciones mediante un conjunto de validación o validación cruzada. Se elige la configuración que tenga mejor desempeño según la métrica definida.
\end{respuesta}

\pregunta{4. ¿Qué conjunto de datos puede utilizarse para comparar diferentes hiperparámetros?}
\begin{respuesta}
El conjunto de \textbf{validación}, o los pliegues de validación en validación cruzada.
\end{respuesta}

\pregunta{5. ¿Por qué no es conveniente seleccionar hiperparámetros utilizando el conjunto de prueba?}
\begin{respuesta}
Porque el test dejaría de ser una evaluación independiente: cada comparación transmite información sobre ese conjunto y puede conducir a elegir una configuración adaptada a sus peculiaridades.
\end{respuesta}

% ============================================================
\section{Interpretación básica de código con scikit-learn}

\begin{lstlisting}
X_train, X_test, y_train, y_test = train_test_split(
    X, y, test_size=0.25, random_state=10
)
model.fit(X_train, y_train)
predictions = model.predict(X_test)
\end{lstlisting}

\pregunta{1. ¿Qué produce \texttt{train\_test\_split}?}
\begin{respuesta}
Divide las entradas $X$ y los objetivos $y$ en subconjuntos correspondientes para entrenamiento y prueba, conservando la alineación entre cada fila de $X$ y su etiqueta $y$.
\end{respuesta}

\pregunta{2. ¿Qué representa \texttt{test\_size=0.25}?}
\begin{respuesta}
Que aproximadamente el 25\% de las observaciones se asignará al conjunto de prueba y el 75\% al entrenamiento.
\end{respuesta}

\pregunta{3. Si tenemos 4,000 observaciones y utilizamos \texttt{test\_size=0.25}, aproximadamente ¿cuántas quedan para prueba?}
\begin{respuesta}
$4000(0.25)=1000$ observaciones.
\end{respuesta}

\pregunta{4. ¿Cuántas quedan para entrenamiento?}
\begin{respuesta}
$4000-1000=3000$ observaciones.
\end{respuesta}

\pregunta{5. ¿Qué ocurre durante \texttt{model.fit(X\_train, y\_train)}?}
\begin{respuesta}
El algoritmo utiliza $X_{train}$ y $y_{train}$ para estimar los parámetros internos del modelo según su procedimiento de entrenamiento.
\end{respuesta}

\pregunta{6. ¿Qué devuelve normalmente \texttt{model.predict(X\_test)}?}
\begin{respuesta}
Un vector de predicciones, una por fila de $X_{test}$. En regresión son valores numéricos; en clasificación normalmente son etiquetas de clase.
\end{respuesta}

\pregunta{7. ¿Por qué hacemos predicciones sobre $X_{test}$ y posteriormente las comparamos con $y_{test}$?}
\begin{respuesta}
Porque $X_{test}$ contiene observaciones no usadas para ajustar el modelo, mientras $y_{test}$ contiene sus respuestas reales. La comparación estima el error fuera de muestra.
\end{respuesta}

\pregunta{8. ¿Qué utilidad tiene establecer \texttt{random\_state}?}
\begin{respuesta}
Fija la semilla del generador pseudoaleatorio usado en la división. Con los mismos datos y el mismo valor se reproduce la misma partición, lo que facilita comparaciones y depuración.
\end{respuesta}

\pregunta{9. ¿Qué diferencia conceptual existe entre \texttt{fit} y \texttt{predict}?}
\begin{respuesta}
\texttt{fit} aprende parámetros a partir de datos etiquetados. \texttt{predict} usa el modelo ya ajustado para producir salidas sobre nuevas entradas.
\end{respuesta}

\pregunta{10. ¿Sería correcto evaluar el modelo comparando \texttt{model.predict(X\_test)} con \texttt{y\_train}? Explica.}
\begin{respuesta}
No. Las predicciones corresponden fila por fila a $X_{test}$ y deben compararse con $y_{test}$. $y_{train}$ pertenece a otras observaciones y normalmente incluso tiene longitud distinta.
\end{respuesta}

% ============================================================
\section{Flujo general de un proyecto de aprendizaje de máquina}

\begin{center}
\begin{tikzpicture}[
  node distance=7mm,
  s/.style={draw,rounded corners,minimum width=9.5cm,minimum height=7mm,fill=blue!4,align=center},
  a/.style={-Latex,thick}
]
\node[s] (a) {1--2. Definir problema y objetivo};
\node[s,below=of a] (b) {3--6. Obtener, explorar, limpiar y construir características};
\node[s,below=of b] (c) {7. Separar datos};
\node[s,below=of c] (d) {8--10. Elegir modelo, entrenar y seleccionar hiperparámetros};
\node[s,below=of d] (e) {11. Evaluar en datos no vistos};
\node[s,below=of e] (f) {12. Desplegar};
\node[s,below=of f] (g) {13--14. Monitorear y reentrenar};
\foreach \u/\v in {a/b,b/c,c/d,d/e,e/f,f/g} \draw[a] (\u)--(\v);
\draw[a,dashed] (g.east) -- ++(1.2,0) |- (b.east);
\end{tikzpicture}
\end{center}

\pregunta{1. ¿Por qué es importante definir primero el problema?}
\begin{respuesta}
Porque determina qué queremos predecir, para quién, con qué datos, bajo qué restricciones y cómo se medirá el éxito. Sin una definición clara es posible optimizar una métrica que no resuelve la necesidad real.
\end{respuesta}

\pregunta{2. ¿Por qué debemos explorar los datos antes de entrenar?}
\begin{respuesta}
Para conocer distribuciones, tipos de variables, escalas, valores faltantes, posibles errores, desbalance de clases, relaciones entre variables y problemas de representatividad.
\end{respuesta}

\pregunta{3. ¿Qué problemas podemos buscar durante la limpieza de datos?}
\begin{respuesta}
Valores faltantes, duplicados, formatos inconsistentes, categorías mal escritas, unidades mezcladas, valores imposibles, registros corruptos y variables que producen fuga de información.
\end{respuesta}

\pregunta{4. ¿Por qué debemos definir una métrica de evaluación?}
\begin{respuesta}
Porque necesitamos una regla cuantitativa que traduzca el objetivo en desempeño medible. La métrica guía la comparación entre modelos y debe reflejar los tipos de errores relevantes.
\end{respuesta}

\pregunta{5. ¿En qué momento se separan los datos de entrenamiento y prueba?}
\begin{respuesta}
Temprano en el flujo, antes de usar información del test para tomar decisiones. Después, transformaciones que aprenden de los datos deben ajustarse sólo con entrenamiento y aplicarse al resto.
\end{respuesta}

\pregunta{6. ¿Qué ocurre durante el entrenamiento?}
\begin{respuesta}
El algoritmo ajusta los parámetros del modelo para reducir una función de pérdida o cumplir un criterio utilizando los datos de entrenamiento.
\end{respuesta}

\pregunta{7. ¿Qué ocurre durante la evaluación?}
\begin{respuesta}
Se generan predicciones sobre datos no usados para ajustar el modelo y se comparan con sus valores reales mediante métricas previamente elegidas.
\end{respuesta}

\pregunta{8. ¿Qué significa desplegar un modelo?}
\begin{respuesta}
Integrarlo en un entorno donde pueda recibir nuevas entradas y producir predicciones para usuarios, sistemas o procesos reales.
\end{respuesta}

\pregunta{9. ¿Por qué un modelo debe monitorearse después de ponerlo en producción?}
\begin{respuesta}
Porque los datos, procesos y poblaciones pueden cambiar. Se deben vigilar desempeño, distribución de entradas, errores, tiempos de respuesta y posibles fallos operativos.
\end{respuesta}

\pregunta{10. ¿Por qué podría disminuir el desempeño de un modelo con el paso del tiempo?}
\begin{respuesta}
Porque puede cambiar la distribución de los datos o la relación entre entradas y objetivo. Este fenómeno suele llamarse \textbf{data drift} o \textbf{concept drift}, según qué parte del proceso cambie.
\end{respuesta}

% ============================================================
\newpage
\section{Ejercicios integradores resueltos}

\subsection{Ejercicio 1: gasto mensual de clientes bancarios}
Un banco quiere estimar cuánto dinero gastará cada cliente con su tarjeta durante el próximo mes.

\begin{enumerate}[label=\alph*)]
\item \textbf{Tipo de aprendizaje:} supervisado, porque se entrenaría con clientes históricos cuyo gasto posterior es conocido.
\item \textbf{Clasificación o regresión:} regresión, porque el gasto es una cantidad numérica continua.
\item \textbf{Variable objetivo:} gasto total del cliente durante el próximo mes.
\item \textbf{Cuatro características posibles:} gasto de los últimos meses, ingreso estimado, número de transacciones recientes y saldo o línea de crédito utilizada. También podrían usarse antigüedad del cliente o estacionalidad.
\item \textbf{Métrica apropiada:} MAE o RMSE. MAE expresa el error típico directamente en unidades monetarias; RMSE penaliza con mayor fuerza errores grandes.
\end{enumerate}

\subsection{Ejercicio 2: condición médica con 150 positivos de 10,000}

\begin{enumerate}[label=\alph*)]
\item \textbf{¿Existe desbalance?} Sí. La clase positiva representa
\[
\frac{150}{10000}=0.015=1.5\%,
\]
y la negativa 98.5\%.
\item Si el modelo siempre predice «no presenta la condición», acierta los $9850$ negativos:
\[
\mathrm{Accuracy}=\frac{9850}{10000}=0.985=98.5\%.
\]
\item Ese 98.5\% no demuestra utilidad porque el modelo tendría $TP=0$ y $FN=150$: no detectaría a ningún paciente positivo, por lo que su recall positivo sería $0$.
\item \textbf{FP:} clasificar como positivo a un paciente que no tiene la condición. \textbf{FN:} clasificar como negativo a un paciente que sí la tiene.
\item Revisaría matriz de confusión, recall, precision y F1-score. Si el objetivo principal es no perder casos positivos, recall merece atención especial.
\end{enumerate}

\subsection{Ejercicio 3: entrenamiento vs. validación}

\begin{center}
\begin{tabular}{lcc}
\toprule
Modelo & Error entrenamiento & Error validación\\
\midrule
M1 & 18\% & 20\%\\
M2 & 8\% & 10\%\\
M3 & 2\% & 17\%\\
\bottomrule
\end{tabular}
\end{center}

\textbf{M1.} Los dos errores son relativamente altos y cercanos. Bajo la interpretación usual de que menor error es mejor, hay evidencia de \textbf{subajuste}: el modelo ni siquiera logra un ajuste fuerte en entrenamiento y tampoco mejora fuera de muestra.

\textbf{M2.} Tiene errores bajos y cercanos: 8\% en entrenamiento y 10\% en validación. Es el patrón más consistente con un modelo que ajusta bien y \textbf{generaliza razonablemente}.

\textbf{M3.} El error de entrenamiento es extremadamente pequeño (2\%) pero el de validación sube a 17\%. La brecha de 15 puntos porcentuales es evidencia clara de \textbf{sobreajuste}.

\begin{advertencia}
El diagnóstico usa sólo la información relativa de la tabla. En una aplicación real también necesitaríamos saber qué significa exactamente ese error y cuál es un nivel aceptable para el problema.
\end{advertencia}

\subsection{Ejercicio 4: métricas de clasificación}
Datos: $TP=40$, $TN=900$, $FP=30$, $FN=30$. El total es
\[
N=40+900+30+30=1000.
\]

\paragraph{a) Accuracy}
\[
\mathrm{Accuracy}=\frac{40+900}{1000}=0.94=94\%.
\]
Interpretación: el 94\% de todas las observaciones fue clasificado correctamente.

\paragraph{b) Precision}
\[
\mathrm{Precision}=\frac{40}{40+30}=\frac{40}{70}\approx0.5714=57.14\%.
\]
Interpretación: de todos los casos que el modelo marcó como positivos, aproximadamente 57.14\% sí eran positivos.

\paragraph{c) Recall}
\[
\mathrm{Recall}=\frac{40}{40+30}=\frac{40}{70}\approx0.5714=57.14\%.
\]
Interpretación: de todos los positivos reales, el modelo detectó aproximadamente 57.14\%.

\paragraph{d) F1-score}
Como precision y recall son iguales,
\[
F_1=2\frac{(0.5714)(0.5714)}{0.5714+0.5714}\approx0.5714=57.14\%.
\]
Interpretación: el balance entre precision y recall es moderado, a pesar de que el accuracy total es alto. Esto ilustra por qué accuracy puede ocultar un desempeño pobre sobre la clase positiva.

\subsection{Ejercicio 5: métricas de regresión}
\[
y=(100,150,200,250),\qquad \hat y=(110,140,180,270).
\]

\paragraph{a) Errores $y_i-\hat y_i$}
\[
(-10,\ 10,\ 20,\ -20).
\]

\paragraph{b) Errores absolutos}
\[
(10,\ 10,\ 20,\ 20).
\]

\paragraph{c) MAE}
\[
\mathrm{MAE}=\frac{10+10+20+20}{4}=15.
\]
El modelo se equivoca, en magnitud absoluta, 15 unidades en promedio.

\paragraph{d) MSE}
\[
\mathrm{MSE}=\frac{100+100+400+400}{4}=250.
\]
Es el promedio de errores cuadrados; castiga más los errores de magnitud 20 que los de magnitud 10.

\paragraph{e) RMSE}
\[
\mathrm{RMSE}=\sqrt{250}=5\sqrt{10}\approx15.811.
\]
Está nuevamente en las unidades originales de $y$ y, debido al cuadrado previo, refleja con mayor sensibilidad los errores grandes que MAE.

\subsection{Ejercicio 6: lectura de código}
\begin{lstlisting}
X_train, X_test, y_train, y_test = train_test_split(
    X, y,
    test_size=0.30,
    random_state=10
)
model.fit(X_train, y_train)
predictions = model.predict(X_test)
\end{lstlisting}

\begin{enumerate}[label=\alph*)]
\item \textbf{Entrenamiento:} aproximadamente 70\% de los datos.
\item \textbf{Prueba:} aproximadamente 30\% de los datos.
\item El modelo aprende usando $X_{train}$ y $y_{train}$.
\item \texttt{predictions} se genera para las observaciones contenidas en $X_{test}$.
\item Las predicciones deben compararse contra $y_{test}$.
\item \texttt{random\_state=10} hace reproducible la parte pseudoaleatoria de la división: al repetir el procedimiento con los mismos datos y parámetros se obtiene la misma partición.
\end{enumerate}

% ============================================================
\newpage
\section{Resumen de memoria rápida}

\begin{center}
\renewcommand{\arraystretch}{1.35}
\begin{tabular}{p{0.27\textwidth}p{0.66\textwidth}}
\toprule
Concepto & Regla mental\\
\midrule
Supervisado & Hay $y$: aprendemos $X\mapsto y$.\\
No supervisado & No hay $y$: buscamos estructura en $X$.\\
Clasificación & Objetivo categórico.\\
Regresión & Objetivo numérico.\\
Train & Aprende parámetros.\\
Validation & Elige modelo e hiperparámetros.\\
Test & Evalúa al final.\\
Subajuste & Error train alto + error no visto alto.\\
Sobreajuste & Error train muy bajo + error no visto bastante mayor.\\
MAE & Magnitud absoluta promedio del error.\\
MSE & Error cuadrático promedio; castiga errores grandes.\\
RMSE & $\sqrt{\mathrm{MSE}}$; mismas unidades que $y$.\\
Precision & De lo predicho positivo, ¿cuánto sí era positivo?\\
Recall & De lo realmente positivo, ¿cuánto encontré?\\
F1 & Balance armónico entre precision y recall.\\
Parámetro & Lo aprende el algoritmo.\\
Hiperparámetro & Lo seleccionamos para configurar el aprendizaje.\\
\bottomrule
\end{tabular}
\end{center}

\begin{clave}
\[
\boxed{\text{Buen ML} = \text{objetivo bien definido} + \text{datos correctos} + \text{validación limpia} + \text{generalización}}
\]
\end{clave}

\end{document}
